% Centri — Formalized evaluation framework (week ending 2026-07-08)
% Reconciles Utami (2025) + P-MAGIC (2026); adds video-recognition + multimodal axes.
% Compile: pdflatex -interaction=nonstopmode ... (run twice). Reuses the 06-30/07-08 preamble.
\documentclass[aspectratio=169,10pt]{beamer}
\usetheme{metropolis}
\metroset{sectionpage=progressbar, subsectionpage=none, numbering=fraction, progressbar=frametitle}

\usepackage{booktabs}
\usepackage{amsmath,amssymb}
\usepackage{siunitx}
\usepackage{xcolor}
\usepackage{colortbl}
\usepackage{graphicx}
\graphicspath{{figs/}}
\usepackage{tikz}
\usetikzlibrary{arrows.meta,positioning,calc,shapes.geometric,fit,backgrounds}
\usepackage[most]{tcolorbox}

% ---- palette (shared with the other Centri decks) ------------------------
\definecolor{cMeas}{HTML}{1F6FB2}
\definecolor{cPed}{HTML}{B5651D}
\definecolor{cApp}{HTML}{2E7D32}
\definecolor{cInfra}{HTML}{5E35B1}
\definecolor{cAccent}{HTML}{C62828}
\definecolor{cGrey}{HTML}{616161}
\definecolor{cBg}{HTML}{F4F6F8}
\definecolor{cBasic}{HTML}{2E7D32}
\definecolor{cInter}{HTML}{1565C0}
\definecolor{cAdv}{HTML}{6A1B9A}
\setbeamercolor{frametitle}{bg=cMeas}
\setbeamercolor{progress bar}{fg=cAccent}

\newcommand{\hl}[1]{{\bfseries\color{cAccent}#1}}
\newcommand{\ok}{{\color{cApp}\textbf{\checkmark}}}

\tikzset{
  box/.style={rounded corners=2pt, draw=#1, line width=0.8pt, fill=#1!8,
              text=black, align=center, font=\scriptsize, inner sep=4pt, minimum height=8mm},
  node distance=5mm and 7mm,
  flow/.style={-{Stealth[length=2.2mm]}, line width=0.7pt, draw=cGrey},
}

\title{Centri --- A Formalized Evaluation Framework}
\subtitle{Reconciling the parent MWP rubric \& P-MAGIC (2026) for video-based, tiered physics material}
\author{Damar Albaribin}
\institute{MSc Thesis --- advisor: Prof.\ Wu-Yuin Hwang}
\date{Week ending 2026-07-08}

\begin{document}
\maketitle

% =========================================================================
\begin{frame}{Two parent rubrics, reconciled for video material}
\small
\begin{tcolorbox}[colback=cGrey!8,colframe=cGrey,title=\footnotesize The two parents,boxsep=2pt,left=4pt,right=4pt]
\scriptsize
The framework reconciles two rubrics from the lab. The \textbf{parent MWP dissertation}
(\textbf{Ika Utami}, NCU, same advisor): its \textbf{SocioMathLLM} study is the LLM-judge $+$ linguistic
rubric (Tables 8--10), its \textbf{Geo-QG} study the AR-sensing / 3-tier precedent. The \textbf{sibling
paper P-MAGIC} (Hwang, Sari et al., \emph{J.\ Educational Computing Research} 2026): same lab, multimodal
physics \emph{word problems} from object-image $+$ phone sensors. Its \textbf{Table 2} is a rubric organised
\emph{by modality} --- the piece needed to evaluate \emph{our} multimodal output.
\end{tcolorbox}
\vspace{0.3em}
Three moves that define the framework:
\begin{itemize}\setlength\itemsep{2pt}
  \item \textbf{\color{cMeas}Reframed for our real unit} \ok\ --- learning \emph{material} (exposition),
  not word problems; the problem-only rubric rows simply don't apply.
  \item \textbf{\color{cPed}Reconciled two parent rubrics} \ok\ --- SocioMathLLM (linguistic) $+$ P-MAGIC
  (multimodal) $+$ Centri (physics/tier $+$ video) into one 4-axis rubric.
  \item \textbf{\color{cApp}Formalized the whole method} \ok\ --- \texttt{eval-framework.md}: axes,
  channels, statistics, and two new contributions.
\end{itemize}
\end{frame}

% =========================================================================
\section{Positioning}

\begin{frame}{The video niche --- and hard differentiation from P-MAGIC}
\scriptsize
P-MAGIC is the \emph{direct sibling}: same advisor, also multimodal, and it \emph{already} does
automated annotation. So multimodality/annotation are \textbf{not} our novelty --- we lead on what differs:
\vspace{0.3em}
\begin{center}
\renewcommand{\arraystretch}{1.3}
\begin{tabular}{@{}lll@{}}
\toprule
 & \textbf{\color{cPed}P-MAGIC (sibling)} & \textbf{\color{cApp}Centri} \\
\midrule
Input & phone IMU bolted to the object & \hl{ordinary video / CV tracking} (nothing attached) \\
Output & multimodal word \emph{problems} & difficulty-tiered learning \emph{material} \\
Grounding & --- & deterministic gate: every number vs ground truth \\
Annotation & one still image $+$ sensor graph & \hl{per-frame, across the motion} \\
\bottomrule
\end{tabular}
\end{center}
\vspace{0.4em}
\begin{tcolorbox}[colback=cAccent!6,colframe=cAccent,boxsep=2pt,left=5pt,right=5pt]
\scriptsize \textbf{The niche (from our own video-lit review):} video analysis is \emph{manual}
(Tracker), CV tracking is a \emph{measurement aid}, VLM question-generation is for \emph{lecture
content} --- \textbf{no one generates grounded, difficulty-tiered physics \emph{material} from
\emph{motion}-video recognition.} (And we are the inverse of the text$\to$video generation crowd.)
\end{tcolorbox}
\end{frame}

% =========================================================================
\section{The formalized method}

\begin{frame}{Four rubric axes, four scoring channels}
\small
Four axes get scored 1--5. Four channels score them --- the \textbf{gate} is ours and runs first
for free; the \textbf{human panel and MLLM-judge share identical rubric rows}.
\vfill
\begin{center}
\renewcommand{\arraystretch}{1.5}
\setlength{\tabcolsep}{9pt}
\begin{tabular}{@{}lcccc@{}}
\toprule
\textbf{Channel} & \textbf{\color{cPed}Ax1}\,{\scriptsize Ling.} & \textbf{\color{cMeas}Ax2}\,{\scriptsize Struct.} & \textbf{\color{cApp}Ax3}\,{\scriptsize Phys+rec} & \textbf{\color{cInfra}Ax4}\,{\scriptsize Multi.} \\
\midrule
\textbf{0 Gate}\, {\scriptsize\color{cGrey}deterministic, ours} & {\color{cGrey}\textendash} & {\color{cGrey}\textendash} & \ok & \ok \\
\textbf{1 Automatic}\, {\scriptsize\color{cGrey}BERT·BLEU·STS} & \ok & \ok & {\color{cGrey}\textendash} & {\color{cGrey}\textendash} \\
\textbf{2 MLLM-judge}\, {\scriptsize\color{cGrey}vision, semi-auto} & \ok & \ok & \ok & \ok \\
\textbf{3 Human}\, {\scriptsize\color{cGrey}teachers, per-modality} & \ok & \ok & \ok & \ok \\
\bottomrule
\end{tabular}\\[5pt]
{\scriptsize\color{cGrey}Axes: \textbf{1} Linguistic \,·\, \textbf{2} Structural / comprehension \,·\,
\textbf{3} Physics-tier + video-recognition \,·\, \textbf{4} Multimodal (image / graph / table)}
\end{center}
\vfill
\begin{center}
\footnotesize\itshape Then correlate across channels: Cohen's $\kappa$ (inter-rater) · ICC
(human$\leftrightarrow$MLLM) · Pearson (auto$\leftrightarrow$human).\\
The gate feeds ground truth into Axes 3--4 --- neither parent has it.
\end{center}
\end{frame}

\begin{frame}{The reconciled rubric --- who each axis comes from}
\scriptsize
\renewcommand{\arraystretch}{1.35}
\begin{center}
\begin{tabular}{@{}p{2.3cm}p{2.6cm}p{3.1cm}p{3.2cm}@{}}
\toprule
\textbf{Axis} & \textbf{\color{cPed}From SocioMathLLM} & \textbf{\color{cMeas}From P-MAGIC} & \textbf{\color{cApp}Centri-only} \\
\midrule
1 Linguistic & Table 8 (all 5) & Linguistic block & --- \\
2 Structural / comprehension & structure \& clarity & concept · realistic · variable-name & comprehension; \emph{drops problem-only rows} \\
3 Physics/tier $+$ recognition & --- & difficulty $\to$ difficulty\_fit & grounding · achieved\_bloom · \hl{video\_recognition\_quality} \\
4 Multimodal & --- & \textbf{image-text · text-graph · text-table} (Table 2) & \hl{annotation\_correctness} (on video) \\
\bottomrule
\end{tabular}
\end{center}
\vspace{0.3em}
{\footnotesize \textbf{Axis 2 is slimmed for exposition:} the parents' problem-only rows
(answerability, \# solution steps, multiple strategies) score a passage as if it were a question, so
they are dropped; the pedagogy signal moves to comprehension / structure / accuracy.\\[3pt]
Adopted rows kept \emph{verbatim} (with P-MAGIC's citations: Otani 2023, Stefanel 2019) so our numbers
sit beside the lab's. Source of truth: \texttt{eval-rubric-ika.md}.}
\end{frame}

% =========================================================================
\section{Contributions \& preliminary signal}

\begin{frame}{Two evaluation contributions the parents structurally cannot make}
\small
\begin{enumerate}\setlength\itemsep{4pt}
  \item \textbf{\color{cApp}Video-recognition quality (Axis 3).} Geo-QG senses via AR, P-MAGIC via a
  phone IMU --- neither recognises from \emph{video}, so neither can score recognition quality. We
  \emph{can}: the grounding gate already computes ground truth, so we score CV-measured
  $\omega, a_c, r, T$ against it (cf.\ the pendulum $r=0.99$ in the CV-tracking precedent).
  \item \textbf{\color{cInfra}MLLM-as-judge over annotated video frames (Axis 4).} A vision model gets
  the \emph{annotated frame $+$ gate ground truth as reference} $\to$ per-criterion 1--5 $+$ rationale
  (the VideoJudge recipe), calibrated on a few teacher-verified items --- the \emph{semi-automatic} channel.
\end{enumerate}
\vspace{0.3em}
\begin{tcolorbox}[colback=cPed!6,colframe=cPed,boxsep=2pt,left=4pt,right=4pt]
\scriptsize \textbf{Annotation quality matters:} P-MAGIC's own ablation shows annotated modalities
beat non-annotated (image$+$annotation $0.896$ vs table$-$annotation $0.822$; Cohen's $d$ up to $1.7$).
Our differentiator is doing it \textbf{on video} --- per-frame, across the motion.
\end{tcolorbox}
\end{frame}

\begin{frame}{Preliminary signal is already there}
\scriptsize
15 passages (5 objects $\times$ 3 tiers), judged by the lab model on the rubric.
\emph{Preliminary --- the teacher panel is next.}
\vspace{0.2em}
\begin{columns}[T]
\begin{column}{0.48\textwidth}
\centering{\scriptsize\bfseries Criteria that should rise with tier}\\[2pt]
\begin{tabular}{@{}lccc@{}}
\toprule
\textbf{Criterion} & \color{cBasic}basic & \color{cInter}int. & \color{cAdv}adv. \\
\midrule
cognitive\_demand & 2.2 & 2.6 & \textbf{4.0} \\
solution\_steps & 1.2 & 2.0 & \textbf{3.0} \\
achieved Bloom & \multicolumn{2}{c}{Understand} & \textbf{Analyze} \\
\bottomrule
\end{tabular}
\end{column}
\begin{column}{0.52\textwidth}
\centering{\scriptsize\bfseries Honest reading}\\[2pt]
\scriptsize
\begin{itemize}\setlength\itemsep{2pt}
  \item Grounding drops at advanced ($4.4\to1.0$) on oblique clips --- now an explicit
  \emph{video\_recognition\_quality} signal, not a footnote.
  \item Automatic baseline holds: BERTScore F1 $=0.830\pm0.014$ (n$=$75).
\end{itemize}
\end{column}
\end{columns}
\vspace{0.4em}
\begin{tcolorbox}[colback=cApp!6,colframe=cApp,boxsep=2pt,left=5pt,right=5pt]
\scriptsize The ladder an independent model sees (cognitive demand $2.2\to4.0$, Bloom
Understand$\to$Analyze) is the pedagogy signal BERTScore can't give --- now inside a formalized,
comparable framework.
\end{tcolorbox}
\end{frame}

% =========================================================================
\begin{frame}{Status \& what's next}
\small
\begin{columns}[T]
\begin{column}{0.52\textwidth}
\textbf{\color{cApp}Done}
\begin{itemize}\setlength\itemsep{1pt}\small
  \item Found $+$ read the P-MAGIC sibling paper \ok
  \item Reconciled 4-axis rubric (SocioMathLLM $+$ P-MAGIC $+$ ours) \ok
  \item Formalized method doc \texttt{eval-framework.md} \ok
  \item Video literature review $+$ hard positioning \ok
  \item New Axis 3 video-recognition $+$ Axis 4 multimodal \ok
  \item Stats/judge/rater tooling already wired \ok
\end{itemize}
\end{column}
\begin{column}{0.48\textwidth}
\textbf{\color{cAccent}Next}
\begin{itemize}\setlength\itemsep{1pt}\small
  \item \textbf{Teacher panel} on the per-modality rubric $\to$ real $\kappa$/ICC/Pearson
  \item \textbf{Wire the MLLM-judge} over annotated frames (mimo-v2.5)
  \item \textbf{Fix video-frame annotation} (arrow$+\omega$, degree overlay, phase figure) --- adopt P-MAGIC's methods
  \item \textbf{Reorder pipeline:} material waits for the annotated video, then generates
  \item Implement the reconciled rubric in \texttt{run\_llm\_judge.py}
\end{itemize}
\end{column}
\end{columns}
\vspace{0.5em}
\begin{center}
{\footnotesize\itshape A formalized, comparable evaluation framework --- reconciled with both parent
lines, extended for video --- ready for the teacher panel to turn preliminary into publishable.}
\end{center}
\end{frame}

\end{document}
